\documentclass{article}
\usepackage{apstats-poster}
% Formula IDs: complement, add-rule, cond-prob, mult-rule, mult-independent
\colorlet{PosterAccent}{UnitBlue}

\begin{document}
\begin{posterpage}
\PosterHeader[88pt]{UnitBlue}{2}{PROBABILITY \& DISTRIBUTIONS}{P05}
  {PROBABILITY RULES}
  {new 2.3--2.7 \quad\textbullet\quad old 4.1--4.6}

\PosterZone{4.72in}{16.08in}{ZONE A \textbullet\ THE FORMULAS}{%
  \begin{minipage}[t]{12.25in}
    \vspace{0pt}
    \FormulaCardSized{46pt}{0\le P(A)\le1\qquad P(S)=1}
    \vspace{.14in}
    \FormulaCardSized{50pt}{P(A^c)=1-P(A)}
    \vspace{.14in}
    \FormulaCardSized{39pt}{\begin{gathered}
      P(A\cup B)=P(A)+P(B)-P(A\cap B)\\[8pt]
      \text{mutually exclusive: }P(A\cap B)=0
    \end{gathered}}
    \vspace{.14in}
    \FormulaCardSized{46pt}{P(A\mid B)=\frac{P(A\cap B)}{P(B)}\quad P(B)>0}
    \vspace{.14in}
    \FormulaCardSized{43pt}{P(A\cap B)=P(A)\cdot P(B\mid A)}
    \vspace{.14in}
    \FormulaCardSized{38pt}{\begin{gathered}
      \text{Independent}\ \Longleftrightarrow\ P(A\mid B)=P(A)\\[8pt]
      \Longleftrightarrow\ P(A\cap B)=P(A)\cdot P(B)
    \end{gathered}}
    \vspace{.1in}
    {\fontsize{25pt}{32pt}\selectfont
      A conditional expression requires a positive probability for its given event.}
  \end{minipage}%
  \hfill
  \begin{minipage}[t]{7.25in}
    \vspace{0pt}
    \MiniHeading{OR: COUNT THE OVERLAP ONCE}\par\vspace{.2in}
    \begin{center}
    \begin{tikzpicture}[x=1in,y=.85in]
      \draw[PosterInk,line width=2.5pt] (.1,.15) rectangle (6.6,4.3);
      \fill[PosterAccent!15] (2.45,2.3) circle (1.52);
      \fill[PosterAccent!15] (4.15,2.3) circle (1.52);
      \begin{scope}
        \clip (2.45,2.3) circle (1.52);
        \fill[PosterAccent!48] (4.15,2.3) circle (1.52);
      \end{scope}
      \draw[PosterAccent,line width=3pt] (2.45,2.3) circle (1.52) (4.15,2.3) circle (1.52);
      \node[font=\bfseries\fontsize{34pt}{37pt}\selectfont] at (1.6,2.3) {$A$};
      \node[font=\bfseries\fontsize{34pt}{37pt}\selectfont] at (5,2.3) {$B$};
      \node[font=\fontsize{27pt}{31pt}\selectfont] at (3.3,2.3) {$A\cap B$};
      \node[anchor=north east,font=\fontsize{26pt}{29pt}\selectfont] at (6.4,4.1) {$S$};
    \end{tikzpicture}
    \end{center}
    \vspace{.2in}
    \MiniHeading{AND: MULTIPLY ALONG A PATH}\par\vspace{.15in}
    \begin{center}
    \begin{tikzpicture}[x=1in,y=.75in]
      \coordinate (start) at (.3,2.4);
      \coordinate (a) at (3,3.6);
      \coordinate (ac) at (3,1.2);
      \draw[PosterAccent,line width=3pt] (start)--(a)--(6,4.4);
      \draw[PosterInk,line width=2.5pt] (a)--(6,2.8) (start)--(ac)--(6,1.7) (ac)--(6,.4);
      \node[anchor=south,rotate=24,font=\fontsize{24pt}{27pt}\selectfont] at (1.5,3) {$P(A)$};
      \node[anchor=north,rotate=-24,font=\fontsize{24pt}{27pt}\selectfont] at (1.5,1.7) {$P(A^c)$};
      \node[anchor=south,fill=white,inner sep=2pt,font=\fontsize{24pt}{27pt}\selectfont] at (4.65,4.05) {$P(B\mid A)$};
      \node[anchor=north,fill=white,inner sep=2pt,font=\fontsize{22pt}{25pt}\selectfont] at (4.65,3.15) {$P(B^c\mid A)$};
      \node[anchor=south,fill=white,inner sep=2pt,font=\fontsize{22pt}{25pt}\selectfont] at (4.65,1.65) {$P(B\mid A^c)$};
      \node[anchor=north,fill=white,inner sep=2pt,font=\fontsize{22pt}{25pt}\selectfont] at (4.65,.65) {$P(B^c\mid A^c)$};
      \node[anchor=west,font=\bfseries\fontsize{25pt}{28pt}\selectfont] at (6.1,4.4) {$B$};
      \node[anchor=west,font=\fontsize{25pt}{28pt}\selectfont] at (6.1,2.8) {$B^c$};
      \node[anchor=west,font=\fontsize{25pt}{28pt}\selectfont] at (6.1,1.7) {$B$};
      \node[anchor=west,font=\fontsize{25pt}{28pt}\selectfont] at (6.1,.4) {$B^c$};
    \end{tikzpicture}
    \end{center}
  \end{minipage}%
}

\PosterZone{16.23in}{20.05in}{ZONE B \textbullet\ WHAT THE SYMBOLS MEAN}{%
  \begin{minipage}[t]{9.7in}
    {\fontsize{32pt}{40pt}\selectfont
      \(\cup\): \textbf{or} \quad \(\cap\): \textbf{and} \quad \(\mid\): \textbf{given}\par
      \(A^c\): event \(A\) does \textbf{not} occur.\par
      \textbf{Outcome:} one result.\par
      \textbf{Event:} a set of outcomes.}
  \end{minipage}%
  \hfill
  \begin{minipage}[t]{9.7in}
    {\fontsize{30pt}{38pt}\selectfont
      \textbf{Sample space \(S\):} all possible outcomes.\par\vspace{.15in}
      \textbf{Law of large numbers:} over many repetitions, an event's
      relative frequency approaches its probability.}
  \end{minipage}%
}

\PosterZone{20.2in}{23.55in}{ZONE C \textbullet\ THE SIMULATION RECIPE}{%
  \begin{minipage}[t]{6.25in}
    \MiniHeading{1 \textbullet\ DESCRIBE ONE TRIAL}\par\vspace{.12in}
    {\fontsize{29pt}{37pt}\selectfont
      Use random digits or the calculator. Match the event and its chance.}
  \end{minipage}%
  \hfill
  \begin{minipage}[t]{6.25in}
    \MiniHeading{2 \textbullet\ REPEAT}\par\vspace{.12in}
    {\fontsize{29pt}{37pt}\selectfont
      Run many trials. Count the trials meeting the success rule.}
  \end{minipage}%
  \hfill
  \begin{minipage}[t]{6.65in}
    \MiniHeading{3 \textbullet\ ESTIMATE}\par\vspace{.12in}
    {\fontsize{30pt}{39pt}\selectfont
      \(\displaystyle P\approx\frac{\#\text{ successes}}{\#\text{ trials}}\)\par\vspace{.1in}
      State the event in context.}
  \end{minipage}%
}

\PosterZone{23.7in}{27.45in}{ZONE D \textbullet\ WATCH OUT}{%
  \begin{minipage}[t]{12.2in}
    \MiniHeading{DISJOINT IS NOT INDEPENDENT}\par\vspace{.12in}
    {\fontsize{33pt}{42pt}\selectfont
      Mutually exclusive events are \textbf{never independent},
      unless at least one has probability 0.}
  \end{minipage}%
  \hfill
  \begin{minipage}[t]{7.3in}
    \MiniHeading{CHECK IT}\par\vspace{.12in}
    {\fontsize{34pt}{43pt}\selectfont
      Independence is checked, not assumed.}
  \end{minipage}%
}
\WorksheetTag{u4\_l1--6}
\end{posterpage}
\end{document}
